\section{ACCOUSTIQUE}

\subsection{Description des manipulations}

L'installation d'un emmeteur d'ondes accoustiques et d'un recepteur piezoelectrique (dont la tension varie en fonction de la pression) sur un banc optique articulé nous ont permi, grâce à un générateur d'impulsions, de mesurer et comparer ces dernières à l'oscilloscope en faisant varier en un premier temps la distance emetteur-recepteur (voir sur Fig. \ref{fig:montage_acc1}). Nous mesurons donc la vitesse de propagation d'un emetteur jusqu'a un recepteur dans un milieu de propagation qui est celui de l'air. Nous mesurons la transmission avec l'emetteur avec des impulsions d'une frequence (f = 40kHz) alimente d'une tension crete a crete ($U_p_p$ = 5V). Cela est possible grace aux mesures et comparer les longueurs d'onde aux deux extremites. Si ces dernieres sont en phase, nous pouvons determiner grace a la distance, la vitesse de propagation de l'onde emise puis recue par le recepteur. Cette distance a egalement un impact sur l'amplitude (ou tension) mesure aux extremites. La tension diminue diminue en fonction de $ln(1/r)$, r etant la distance entre l'emetteur et le recepteur. 

\begin{figure}[ht]
    \centering
    \includegraphics[scale=0.4]{montage_acc1.png}
	\caption{Montage experience accoustique (Vue de profil)}
        \label{fig:montage_acc1}
\end{figure}

En un second temps, nous avons pu determiner la nature spherique d'une onde accoustique dans l'air. En effet, pour ce faire, nous deplacons l'emetteur E, fixe sur le banc articule (voir sur Fig. \ref{fig:montage_acc2}), et mesurons les differentes tensions du recepteur R. 
\begin{figure}[ht]
    \centering
    \includegraphics[scale=0.6]{montage_acc2.png}
	\caption{Montage experience accoustique (Vue de haut)}
        \label{fig:montage_acc2}
\end{figure}
